\documentclass{article}

\usepackage{tabularx}
\usepackage{booktabs}

\title{Problem Statement and Goals\\\progname}

\author{\authname}

\date{}

\input{../Comments.text}
\input{../Common.text}

\begin{document}

\maketitle

\begin{table}[hp]
\caption{Revision History} \label{TblRevisionHistory}
\begin{tabularx}{\textwidth}{llX}
\toprule
\textbf{Date} & \textbf{Developer(s)} & \textbf{Change}\\
\midrule
Date1 & Name(s) & Description of changes\\
Date2 & Name(s) & Description of changes\\
... & ... & ...\\
\bottomrule
\end{tabularx}
\end{table}

\section{Problem Statement}

\wss{You should check your problem statement with the
\href{https://smiths.github.io/capTemplate/Checklists/ProbState-Checklist.pdf}
{problem statement checklist}}. 

\wss{You can change the section headings, as long as you include the required
information.}

\subsection{Problem}

\wss{What problem are you trying to solve?  Why is it important? Do not talk
about how you will solve the problem, only what the problem is.  AI is rarely
part of the problem.  The problem is ``what'' you want to do, not ``how'' you
want to do it.  The problem statement should be written in a way that is
understandable to a non-technical audience.}

\subsection{Inputs and Outputs}

\wss{Characterize the problem in terms of ``high level'' inputs and outputs.  
Use abstraction so that you can avoid details.}

\subsection{Stakeholders}

\subsection{Environment}

\wss{Hardware and Software Environment for the users.  The developer environment
is summarized as part of the developer plan.}

\section{Goals}

This project has a total of five primary goals.

\begin{itemize}
    \item[2.1] \textbf{Automated Decentralized Execution of Transactions and Agreements:}
    Develop a system that enables the automated and decentralized execution of marketplace transactions and agreements between users and model providers.

    \item[2.2] \textbf{Global AI Model Marketplace:}
    Develop a global marketplace that connects users with AI models, allowing models to be discovered and ranked based on their reputation.

    \item[2.3] \textbf{Linear Regression Inference Support:}
    Support linear regression inference tasks, allowing users to submit inference requests to available models through the decentralized marketplace.

    \item[2.4] \textbf{Consumer-Side User Interface:}
    Develop an intuitive consumer-side interface that allows users to interact with the marketplace and submit inference requests without requiring knowledge of the underlying decentralized infrastructure.

    \item[2.5] \textbf{Decentralized Backend Communication for Model Providers:}
    Enable model providers to communicate with and interact with the decentralized backend in order to register their models, receive inference requests, and provide inference results.
\end{itemize}


\section{Stretch Goals}
This project has a total of 2 stretch goals.

\begin{itemize}
    \item[3.1] \textbf{Support for More Complex Inference Tasks:}
    Extend the system beyond linear regression to support more complex inference tasks and AI models.

    \item[3.2] \textbf{Private Inference Using Trusted Execution Environments:}
    Investigate and implement Trusted Execution Environments (TEEs) to enable inference tasks to be executed on completely private data while protecting the data from unauthorized access.
\end{itemize}


\wss{Goals you hope to get to, but not necessary.  They are also helpful if the
scope of the original project needs to be expanded.}

\section{Custom Documents (CDs)}

\wss{For CAS 741: State whether the project is a research project. This
designation, with the approval (or request) of the instructor, can be modified
over the course of the term.}

\wss{For SE Capstone: List your custom documents (CDs).  Potential CDs include
usability testing, code walkthroughs, user documentation, formal proof,
GenderMag personas, Design Thinking, etc.  (The full list is on the course
outline and in Lecture 02.) Normally the number of CDs will be two (one per
term).  Approval of the CDs will be part of the discussion with the instructor
for approving the project. The CDs, with the approval (or request) of the
instructor, can be modified over the course of the term.}

\newpage{}

\section*{Appendix --- Reflection}

\wss{Not required for CAS 741}

\input{../Reflection.text}

\begin{enumerate}
    \item What went well while writing this deliverable? 
    \item What pain points did you experience during this deliverable, and how
    did you resolve them?
    \item How did you and your team adjust the scope of your goals to ensure
    they are suitable for a Capstone project (not overly ambitious but also of
    appropriate complexity for a senior design project)?
\end{enumerate}  

\end{document}